\documentclass[12pt]{article}

\usepackage[margin=1in]{geometry}
\usepackage{mathptmx}
\usepackage{setspace}
\usepackage{microtype}
\usepackage{amsmath,amssymb}
\usepackage{booktabs}
\usepackage{graphicx}
\usepackage{longtable}
\usepackage[super,sort&compress]{natbib}
\usepackage[hidelinks]{hyperref}
\usepackage{xcolor}

\onehalfspacing

\newcommand{\result}[1]{\textbf{\textcolor{blue}{[RESULT: #1]}}}
\newcommand{\methodtodo}[1]{\textbf{\textcolor{purple}{[METHOD: #1]}}}
\newcommand{\datatodo}[1]{\textbf{\textcolor{teal}{[DATA: #1]}}}

\renewcommand{\thesection}{Supplementary Note \arabic{section}}
\renewcommand{\thetable}{S\arabic{table}}
\renewcommand{\thefigure}{S\arabic{figure}}

\title{\textbf{Supplementary Information}\\[4pt]
\large A molecular fragmentation world reveals the organization of the dark metabolome}
\author{}
\date{}

\begin{document}
\maketitle

\tableofcontents
\newpage

\section{Scope and reading order}

This document contains the supporting experiments for two coupled contributions: a shared molecular fragmentation World supporting Free, Guided and Inverse inference, and the use of the resulting evidence hierarchy to resolve and analyse dark chemistry at scale. The notes follow evidentiary dependency rather than main-figure order. World semantics, evidence levels and biological family definitions are frozen before the downstream analyses they support.

Every quantity marked \result{} must be bound to an immutable artifact recorded in the claim--evidence matrix. No value is transferred from a training log or an exploratory analysis.

\section{Corpus construction, snapshots and leakage control}

\subsection{Library snapshot reconstruction}
\methodtodo{Reconstruction of GNPS, MassBank and MoNA at the freeze and evaluation dates; per-structure first-appearance date; depositing accession and laboratory; handling of re-deposited and revised entries; treatment of structures whose first appearance cannot be dated.}

Structures whose deposition date cannot be established are excluded from the prospective set rather than assumed recent. Table~\ref{tab:snapshots} reports snapshot composition.

\begin{table}[h]
\centering
\caption{Library snapshot composition at freeze and evaluation dates.}
\label{tab:snapshots}
\begin{tabular}{lrrr}
\toprule
Source & Structures at freeze & Structures at evaluation & New in interval \\
\midrule
GNPS & \result{n} & \result{n} & \result{n} \\
MassBank & \result{n} & \result{n} & \result{n} \\
MoNA & \result{n} & \result{n} & \result{n} \\
Union (deduplicated) & \result{n} & \result{n} & \result{n} \\
\bottomrule
\end{tabular}
\end{table}

\subsection{Spectral corpora}
\methodtodo{Provenance, spectrum counts, adduct, instrument and collision-energy distributions for the repository corpus, the mouse cohort, the benchmark corpora and the calibration split.}

\subsection{Definition of the dark subset}
\methodtodo{The frozen library-search protocol: search engine and version, tolerances, scoring, acceptance threshold, and the validation showing that the threshold does not admit false library hits at a rate that would contaminate the dark subset.}

The dark subset is defined by absence of an accepted spectral-library match under this protocol. Dark does not mean the structure is absent from every chemical database; a molecule may exist in PubChem or HMDB while its experimental spectrum remains unannotated. This distinction is preserved throughout.

\subsection{Leakage audit}
\methodtodo{Exact-structure and close-analogue overlap between every evaluation set and every pretraining corpus, including the molecular corpora used for METEOR pretraining and any simulator-generated training data. Analogue exclusion at Tanimoto $\geq 0.7$ or MCES $< 10$.}

Simulator-generated training data requires specific attention: pseudo-spectra generated from structures that appear in an evaluation set constitute leakage even though no experimental spectrum was shared.

\section{Decision experiment 1: constructibility of the prospective set}

\textbf{Purpose.} Establish whether a blinded prospective assessment can be built from library growth.

\textbf{Procedure.} Intersect structures deposited between the freeze and evaluation snapshots with spectra present but unannotated in the repository corpus at freeze. Exclude any structure whose close analogue existed at freeze. Audit depositing laboratories for independence from this work.

\textbf{Pass criterion.} At least 300 qualifying spectra spanning at least three chemical superclasses and five depositing groups.

\textbf{Consequence of failure.} Calibration reverts to held-out library spectra with a scaffold-disjoint stress test, and the atlas is presented as a large-scale application rather than as a released resource. This changes the manuscript and must be reported to all authors before proceeding.

\result{Insert qualifying spectrum count, superclass and depositing-group coverage, and exclusion accounting.}

\section{Decision experiment 2: the soundness--resolution trade-off}

\textbf{Purpose.} Establish whether mechanistic elimination can be operated at a controlled false-exclusion rate while still resolving a useful fraction of spectra.

\textbf{Procedure.} On leakage-controlled benchmark spectra with known answers and matched candidate pools, sweep elimination stringency. At each setting record the rate at which the true structure is eliminated and the resulting distribution over resolution levels.

\textbf{Pass criterion.} At a false-exclusion rate of at most 5\%, at least 40\% of spectra resolve to a unique structure, a bounded isomer set, or a chemical class.

\textbf{Consequence of failure.} The current bounded-statement measurement is not reliable enough to underwrite the census. Figure 1 is redesigned around calibrated candidate sets rather than elimination, and downstream entity counting is repeated under that output definition.

\result{Insert the full stringency sweep, the selected operating point and its justification.}

\section{The molecular fragmentation World}

\subsection{State, transitions and observation semantics}
\methodtodo{Canonical WGV release; bond-electron state; explicit atom/H identity; formal charge and radicals; exact mass; ionization microstates; charged and neutral coproducts; sparse electron-flow actions; multi-branch siblings; branch-local STOP; reconvergent path handling; observation model; engineering censoring and residual frontier mass.}

\subsection{Chemical legality and learned preference}
\methodtodo{Hard constraints for atom/electron/charge/H/mass/formula/valence conservation and materialization; definition of legal action masks; learned transition propensity and branch probability; distinction between executable chemistry and validated mechanism.}

\subsection{Support, specificity and elimination}
For a candidate $M$ and observed peak $p$, the model returns whether some atom-balanced trajectory of $M$ reaches a fragment whose exact mass matches $p$ within tolerance. Three distinct quantities follow and are never conflated:

\begin{enumerate}
\item \textbf{support} --- such a trajectory exists;
\item \textbf{specificity} --- the support exceeds the best support available to competing candidates in the same pool;
\item \textbf{elimination} --- no trajectory of $M$ reaches $p$, so $M$ is excluded at the operating stringency.
\end{enumerate}

Elimination is the operation that makes the output composable. It is also the operation that can be wrong, because the reaction basis is incomplete: absence of a trajectory is evidence of impossibility only to the extent that the model covers the relevant chemistry. Supplementary Note 4 quantifies this directly, and the operating stringency is chosen from that measurement rather than assumed.

\subsection{Accidental-match opportunity control}
\methodtodo{Correction for the probability that a candidate supports a peak by coincidence, as a function of support-set size, mass tolerance and peak density. Without this control, candidates with large support sets are favoured for reasons unrelated to identity.}

\subsection{Trajectory authority}
\methodtodo{Levels of mechanistic authority: mass-compatible support, atom-balanced trajectory, and trajectory consistent with an established gas-phase mechanism. Only the highest level supports a reaction-level claim in the main text.}

\subsection{Computational cost}
\result{Insert per-candidate enumeration time, memory, scaling with heavy-atom count, and total cost of the atlas run.}

\section{World-based structural evidence hierarchy}

\subsection{Definition}
Each spectrum or candidate may reach progressively stronger evidence:
\begin{enumerate}
\item precursor mass / molecular-formula compatibility;
\item chemically valid molecular structure;
\item reachability in the learned fragmentation World;
\item executable reaction-path support for observed evidence;
\item candidate-specific discriminative evidence;
\item bidirectional spectrum--structure consistency under World replay;
\item calibrated bounded structural statement;
\item authentic-standard confirmation.
\end{enumerate}

These levels are frozen before repository-scale deployment. Stronger evidence is expected to increase reliability while reducing coverage; both quantities are reported.

\subsection{Bounded structural output}
At the calibrated operating point the reported structural statement may be a unique putative structure, bounded isomer set, shared substructure or class, molecular formula only, or unresolved. Only authentic-standard evidence is termed identification.

\subsection{Relation to established reporting conventions}
\methodtodo{Mapping onto the community confidence-level scheme, and explicit statement of the difference: the community levels are assigned by human judgement after annotation, whereas these levels are computed from which candidates the measurement eliminates. The two are reported side by side so that readers can translate.}

\subsection{Confidence model and calibration}
\methodtodo{Definition of the confidence that the reported statement contains the true structure; features used; fitting procedure; the disjoint calibration split; the estimator used for calibration error and its confidence intervals.}

The calibration mapping is fitted only on the calibration split and is frozen before the prospective evaluation. It is never refitted after seeing prospective outcomes.

\section{Prospective validation}

\subsection{Freeze protocol}
\methodtodo{Artifacts frozen and hashed before prospective evaluation: model weights, fragmentation rule set, candidate-generation configuration, calibration mapping, evaluation code. Verification that the frozen artifacts reproduce the calibration-split results exactly.}

\subsection{Independence audit}
\methodtodo{Verification that depositing laboratories are independent of this work; exclusion of any structure whose determination could have been influenced by our predictions; treatment of structures deposited by consortia in which authors participate.}

\subsection{Results}
\result{Insert calibration error, containment at threshold, stratification by structural distance from the frozen library, by chemical class, by instrument and by collision energy.}

\subsection{Comparator calibration}
\methodtodo{Comparator scores mapped to confidences by the same procedure and evaluated on the identical prospective set. Comparators are given their best available calibration, not a naive score-to-probability transform.}

\result{Insert comparator calibration errors and the comparison to mechanistic elimination.}

\section{Free, Guided and Inverse inference}

\subsection{Shared World representation}
\methodtodo{Shared graph/frontier/spectrum encoders; task-mode conditioning; checkpoint lineage; state dimensions; target-spectrum firewall for Free; shared versus task-specific heads.}

\subsection{Free inference}
\methodtodo{Subtrajectory-balance/GFlowNet objective; multibranch probability flow; sibling/STOP normalization; observation likelihood; training data and censoring handling.}

\subsection{Guided inference}
\methodtodo{Spectrum-conditioned frontier/value model; DAgger on current-policy states; goal-conditioned IQL; matched-budget planning; explanatory-path endpoints and no-chemistry-change constraint relative to Free.}

\subsection{Inverse inference}
\methodtodo{Heavy-atom construction state; canonical graph identity; formula/charge/H/valence masks; conditional GFlowNet/SMC; posterior diversity; forward World replay and bidirectional consistency.}

\subsection{Optional language-model reasoning interface}
\methodtodo{If METEOR language-model traces are retained in the final system, define their role only as a candidate/hypothesis interface to the canonical WGV system; specify which trace claims are checked against the World and which are not used for online chemistry decisions.}

Language-model trace quality is a supporting diagnostic. Main scientific claims are based on held-out World/Guided/Inverse endpoints and downstream evidence, not token-level reasoning accuracy.

\section{Outer-loop scientific program evolution}

\subsection{Failure localization}
\methodtodo{How systematic World/Guided/Inverse failures are assigned to state representation, action language, observation model, search or training objective; receipt format and failure clustering.}

\subsection{Typed revisions and acceptance gates}
\methodtodo{Allowed program/state/action/algorithm changes; deterministic chemistry compilation; local causal replay; ablation; cross-family transfer; molecule-disjoint held-out gate; regression checks.}

Low-cost worker models or LLM agents may propose revisions. They cannot accept them. Promotion requires frozen deterministic/scientific evaluation.

\subsection{Results}
\result{Insert proposed/rejected/promoted revision counts, scientific KPI changes by revision class and compute cost.}

\section{Repository-scale dark-matter resolution}

\subsection{Evidence-level assignment}
\methodtodo{Frozen thresholds for all World-based evidence levels; handling of unresolved and resource-censored cases; per-spectrum evidence receipts; strongest-evidence assignment rule.}

\subsection{Entity collapsing}
\methodtodo{Rules for collapsing adducts, in-source fragments, isotopologues, charge states and redundant acquisitions onto a single molecular entity; validation of these rules on known compounds with known ion forms; the residual error rate of collapsing and its effect on the entity count.}

The entity count is the headline number of Figure 2 and is only as good as this validation. Report the count with an uncertainty interval derived from the measured collapsing error rate, not as a point estimate.

\subsection{Recurrence}
\methodtodo{Criteria for an entity to be counted as recurrent; independence of datasets; treatment of datasets sharing samples or laboratories.}

\subsection{Per-entry record schema}
\methodtodo{Fields released for every atlas entry: representative spectrum, surviving candidate set, eliminating peaks, supporting trajectories, resolution level, calibrated confidence, recurrence, source datasets, and for unresolved entries the measurement that would resolve them.}

\section{Structural organization of high-evidence dark chemistry}

\subsection{Distance to known metabolite space}
\methodtodo{The frozen reference metabolite set; the preregistered transformation vocabulary; how distance is computed for entities resolved below unique-structure level; sensitivity to the choice of reference set.}

Entities resolved only to class level cannot be assigned a transformation distance and are reported separately rather than imputed.

\subsection{Chemical class composition}
\methodtodo{Class assignment procedure; comparison against reference library composition; identification of classes absent from library coverage; correction for ionization and detectability bias, which affects what enters the corpus at all.}

\subsection{Transformation series}
\methodtodo{Construction of candidate transformation series from structural relationships; the explicit rule that a mass difference plus a structural relationship does not establish an enzymatic reaction; what would be required to make that claim.}

\section{Biological origin attribution}

\subsection{Dataset assembly}
\methodtodo{Query over repository metadata for mouse datasets containing germ-free or gnotobiotic controls, antibiotic treatment or defined dietary intervention; manual verification of every included accession against its publication; exclusion criteria; final inventory with sample counts and platforms.}

Automated metadata parsing identifies candidates only. Every included dataset is manually verified against its associated publication, because repository metadata frequently misdescribes experimental design.

\subsection{Presence/absence modelling}
\methodtodo{Detection model per entity per sample; batch, platform and laboratory controls; handling of missing-not-at-random detection; the statistical model producing an attribution.}

\subsection{Positive controls and nulls}
\methodtodo{Panel of metabolites with established microbial, dietary or host origin; required recovery rate before any dark attribution is interpreted; label-permutation null.}

\result{Insert positive-control recovery rate and permutation null distribution.}

\subsection{Boundary of the claim}
Attribution establishes dependence on a perturbation, not biosynthetic route. Manuscript language uses microbiota-dependent, diet-dependent and host-associated throughout, and families consistent with more than one origin are reported as unresolved rather than assigned to the most likely.

\section{Feature-versus-family biological analysis}

\subsection{Freezing}
\methodtodo{Structural-family definitions and biological contrasts hashed before any association testing; the record demonstrating that family membership and the primary global comparison endpoint were frozen before phenotype analysis.}

\subsection{Paired feature-versus-family analysis}
\methodtodo{Run identical samples, preprocessing and covariates in two arms: anonymous dark features and frozen structural families. Pre-register one primary global endpoint, preferably cross-cohort replication rate or effect-sign concordance. Report the paired difference over the complete eligible family/contrast universe rather than selected positives.}

\result{Insert the primary feature-versus-family global effect with confidence interval and denominator.}

\subsection{Principal programme and replication}
\methodtodo{Selection rule for the principal family after the global comparison is frozen; biological replicate as unit; covariates; multiplicity control; independent replication; leave-one-cohort-out criterion.}

\subsection{Sensitivity}
\methodtodo{Global and principal-family results as a function of confidence/resolution thresholds; metadata permutation; family-definition perturbation; cohort exclusion.}

A family-level biological claim that appears only after changing family membership or confidence thresholds is reported as exploratory rather than confirmatory.

\section{Orthogonal structure validation}

\subsection{Anchor selection}
\methodtodo{The frozen discovery criterion used to select anchors before any standard is purchased; restriction to commercially available compounds; the record showing selection preceded purchase.}

\subsection{Acceptance criteria}
\methodtodo{Precursor mass, MS/MS similarity and retention/coelution criteria for standard-confirmed identification; the chromatographic conditions; treatment of diastereomers and other alternatives that co-elute.}

\subsection{Results including failures}
\result{Insert per-anchor outcome, including anchors that failed to confirm and the reason.}

Anchors that fail to confirm are reported. An unconfirmed anchor is informative about the system's error modes and its omission would misrepresent the validation rate.

\subsection{Unresolved cases}
\methodtodo{For at least one representative case: the surviving isomer set, the evidence that eliminated all other candidates, why the survivors cannot be separated by the available data, and the specific measurement that would separate them.}

\section{Limitations}

\begin{enumerate}
\item \textbf{Mechanistic completeness is empirical.} Elimination stringency is calibrated to a measured false-exclusion rate. Chemistry poorly covered by the fragmentation model is protected by lower reported confidence, not by a guarantee.
\item \textbf{The prospective set is not a random sample.} It contains chemistry that independent laboratories chose to characterize during the freeze interval, which over-represents tractable and biologically topical compounds.
\item \textbf{Origin attribution establishes dependence.} No biosynthetic experiment was performed here.
\item \textbf{Detectability bias.} The atlas describes molecules that ionize and fragment under the acquisition conditions represented in the corpus. It is not a census of the metabolome.
\item \textbf{Structures above the reference-standard level are hypotheses.} They are reported with calibrated confidence and are not identifications.
\item \textbf{No new measurement was acquired.} Evidence-guided acquisition, in which the system selects the next experiment and that experiment is performed, is described as a future direction and is not demonstrated.
\end{enumerate}

\section{Supplementary tables}

\methodtodo{S1 corpus inventory; S2 leakage audit; S3 prospective set composition and independence audit; S4 stringency sweep; S5 comparator configurations and versions; S6 rule-evolution receipts; S7 atlas entry schema; S8 uncharacterized recurrent families; S9 perturbation dataset inventory with manual verification status; S10 anchor dossiers.}

\bibliographystyle{unsrtnat}
\bibliography{../references}

\end{document}
